\section{Strategy 2: Diagonal Discharging}
\label{sec:strategy2-discharging}

Strategy~2 is most naturally read as a discharging argument on the hexagon
skeleton.  A prescribed diagonal demand reduces the boundary capacity of its
V triangle.  The resulting loss is discharged to the next boundary edge.  A
center interval may absorb part of this charge, while a center-free
nonsupercritical V triangle transports the remaining charge without loss.  The
exact \(g\)-maps are retained: they are the local constitutive laws that say
how much charge a given diagonal demand releases at a given incoming boundary
state.

This viewpoint is not a second proof layered on top of the transfer proof.
It is an additive form of the same exact inequalities.  It separates the
proof into three levels:
\begin{enumerate}
\item local geometry determines a release function;
\item the skeleton transports and combines released charge;
\item a terminal capacity, threshold, or radial-separation inequality yields
      the contradiction.
\end{enumerate}

\subsection{The exact local discharge rule}

For \(0\le a,c\le1\), write
\[
 F_c(a):=\widehat g_c(1-a),
 \qquad
 G_c(a):=\widehat g_c^\vee(a)=1-F_c(a).
\]
Thus \(F_c(a)\) is the exact proof-safe outgoing cap for a
nonsupercritical V triangle with incoming demand at least \(a\) and own-radial
demand at least \(c\), while \(G_c(a)\) is the following incoming demand on a
center-free edge.

\begin{definition}[Diagonal release and center subsidy]
\label{def:body-diagonal-release}
Define the \emph{diagonal release}
\begin{equation}
 \mathfrak d_c(a)
 :=G_c(a)-a
 =1-a-F_c(a)
 \ge0.
 \label{eq:body-diagonal-release}
\end{equation}
For an empty or closed center interval \(J\subseteq[0,1]\), let
\(e_J(p)\) be the connected covered extent from the left endpoint, as in
Lemma~\ref{lem:reader-generalized-handoff}, and define the
\emph{center subsidy}
\begin{equation}
 \mathfrak t_J(p):=e_J(p)-p\ge0.
 \label{eq:body-center-subsidy}
\end{equation}
\end{definition}

\begin{proposition}[Exact discharging identity]
\label{prop:body-exact-discharging}
Let \(T\) be a nonsupercritical actual V triangle with
\(A(T)\ge a\) and \(C(T)\ge c\).  Let \(J\) be the center interval on its
outgoing edge, and suppose the outgoing trace of \(T\), the interval \(J\),
and the incoming trace of the next incident V triangle cover that edge.  Then
\begin{equation}
 A(T_{\rm next})
 \ge
 a+\mathfrak d_c(a)
   -\mathfrak t_J\!\left(F_c(a)\right).
 \label{eq:body-exact-discharging}
\end{equation}
For \(J=\varnothing\), this becomes
\begin{equation}
 A(T_{\rm next})\ge a+\mathfrak d_c(a)=G_c(a).
 \label{eq:body-center-free-discharge}
\end{equation}
\end{proposition}

\begin{proof}
The generalized handoff gives
\[
 A(T_{\rm next})
 \ge \mathcal R_J(F_c(a))
 =1-e_J(F_c(a)).
\]
By Definition~\ref{def:body-diagonal-release},
\[
 1-e_J(F_c(a))
 =1-F_c(a)-\bigl(e_J(F_c(a))-F_c(a)\bigr)
 =a+\mathfrak d_c(a)-\mathfrak t_J(F_c(a)).
\]
For the empty interval, \(e_\varnothing(p)=p\).
\end{proof}

If \((A,B,C)\) are the actual reaches of a nonsupercritical V triangle, its
actual unused boundary capacity is
\[
 s(T):=1-A-B\ge0.
\]
Whenever \(C\ge c\), the exact outgoing cap gives
\begin{equation}
 s(T)\ge1-A-F_c(A)=\mathfrak d_c(A).
 \label{eq:body-slack-dominates-release}
\end{equation}
Thus \(\mathfrak d_c(A)\) is the portion of the actual capacity slack that is
forced solely by the diagonal demand and the current incoming state.

\subsection{The discharging dictionary}

The transfer alphabet used below has the following additive interpretation.
\begin{table}[H]
\centering
\scriptsize
\renewcommand{\arraystretch}{1.12}
\begin{tabular}{@{}p{.24\textwidth}p{.28\textwidth}p{.37\textwidth}@{}}
\toprule
exact or relaxed transfer&discharging operation&mathematical meaning\\
\midrule
\(\mathrm I(x)=x\)&wire&transport the current boundary charge unchanged\\
\(G_c(x)\)&exact diagonal release&deposit \(\mathfrak d_c(x)=G_c(x)-x\)\\
\(x+\lambda(x-d)\)&affine deposit&deposit the certified amount
\(\lambda(x-d)\) on its stated selected-\(T_+\) domain\\
threshold at \(e(d)\)&gated reservoir&after an input crosses \(e(d)\), raise
it to at least \(1-e(d)\)\\
\(\mathcal R_J\)&center subsidy&subtract only the connected extension
\(\mathfrak t_J\), not the full length of \(J\)\\
\(g_c^{\rm sc}\)&supercritical sink cap&a strict-supercritical V triangle
with radial demand \(c\) sends out strictly less than \(g_c^{\rm sc}\)\\
path budget&capacity audit&the internal nonsupercritical V triangles can
absorb at most one unit each\\
\bottomrule
\end{tabular}
\caption{The additive interpretation of the Strategy~2 transfer alphabet.}
\label{tab:body-discharging-dictionary}
\end{table}

For a center-free formal chain
\[
 z_j=G_{c_j}(z_{j-1})
 \qquad(1\le j\le r),
\]
there is an exact telescoping identity
\begin{equation}
 z_r-z_0
 =\sum_{j=1}^r\mathfrak d_{c_j}(z_{j-1}).
 \label{eq:body-discharge-telescope}
\end{equation}
The dependence on the arriving state is essential.  A diagonal can release
zero charge on a Full face and a positive amount after earlier roles have
raised the incoming state.

\subsection{Center-free paths and paired endpoint release}

Let \(T_1,\ldots,T_k\) be a center-free ordinary boundary path with external
covered extents \(u\) and \(v\) at its two ends.  Lemma
\ref{lem:reader-boundary-path} gives
\[
 \sum_{j=1}^k(A_j+B_j)\ge k+1-u-v.
\]
If all path roles are nonsupercritical, their total capacity is at most
\(k\).  Hence a positive endpoint surplus
\begin{equation}
 1-u-v>0
 \label{eq:body-positive-endpoint-surplus}
\end{equation}
is enough for a contradiction.

The endpoint certificates have an exact charge form.  Suppose two endpoint
roles have incoming demands \(a_L,a_R\), diagonal demands \(c_L,c_R\), and
outgoing caps
\[
 u_L=F_{c_L}(a_L),
 \qquad
 u_R=F_{c_R}(a_R).
\]
Then
\begin{equation}
 u_L+u_R<1
 \quad\Longleftrightarrow\quad
 \mathfrak d_{c_L}(a_L)+\mathfrak d_{c_R}(a_R)
 >1-a_L-a_R.
 \label{eq:body-endpoint-release-equivalence}
\end{equation}
Thus the two endpoint diagonals jointly release more charge than the endpoint
tax \(1-a_L-a_R\).  The positive remainder is discharged into the ordinary
middle path.

\subsection{The all-\(\Vdzero\) gap-rank kernel}

Let \(\mathrm{gr}\) be the number of positive center traces containing a
V-gap.  The ranks \(0,1,2\) are exhaustive, with rank two possible only for
\(\CEtwo\).

\paragraph{Rank zero with \(N_+=0\).}
No center trace is needed on any boundary handoff.  Since the incident roles
are open, each edge gives
\[
 B_i+A_{i+1}>1.
\]
Summing cyclically,
\[
 \sum_{i=0}^5(A_i+B_i)
 =\sum_{i=0}^5(B_i+A_{i+1})>6.
\]
But all six V triangles are nonsupercritical, so the left side is at most
\(6\), a contradiction.  This is the summed discharging form of the former
strict cyclic ascent.

\paragraph{Rank one.}
The center subsidy is confined to one endpoint edge.  The one-side endpoint
theorem proves
\[
 F_{c_5}(a_5)+F_{c_1}(a_1)<1.
\]
By \eqref{eq:body-endpoint-release-equivalence}, the two forced diagonal
releases exceed the endpoint tax.  Their positive remainder enters the three
middle nonsupercritical V triangles, whose total capacity is only three.

\paragraph{Rank two.}
The common \(\CEtwo\) paired-endpoint theorem gives the same strict endpoint
sum with both center traces retained.  Again, the two endpoint releases
jointly exceed the endpoint tax, and the center-free three-V-triangle path
cannot absorb the remainder.

\subsection{The one-gap state with one supercritical V triangle}

Use the signed center variables
\[
 X=R-\delta,
 \qquad
 H=\frac{\eta+\alpha+\delta}{2R},
 \qquad
 m_3=\min\left\{\frac\alpha R,\frac\delta W\right\},
\]
and the five nonsupercritical radial demands
\[
 \begin{aligned}
 c_1&=1-\frac\delta R,&c_2&=1-\delta,&c_3&=1-m_3,\\
 c_4&=1-\alpha,&c_5&=1-\frac\alpha W.
 \end{aligned}
\]
The exact five-V-triangle discharge starts with charge \(X\) and returns
\begin{equation}
 Z=[G_{c_1}\mid G_{c_2}\mid G_{c_3}\mid G_{c_4}\mid G_{c_5}](X).
 \label{eq:body-five-discharge-chain}
\end{equation}
The actual-V-triangle interface gives
\[
 A_0\ge Z,
 \qquad
 A_0<1-H.
\]
Thus the terminal target is \(Z>1-H\).  For \(\CEone\), extensivity of
the first and fifth slots and exact duality reduce this target to the reversed
three-reservoir problem
\begin{equation}
 [G_{1-\alpha}\mid G_{1-m_3}\mid G_{1-\delta}](H)>1-X.
 \label{eq:body-three-discharge-chain}
\end{equation}
For \(\CEtwo\), the threshold dichotomy is more naturally used in the
original five-V-triangle order, starting from \(X\) and targeting
\(1-H\).

\paragraph{The \(\CEone\) deposit cascade.}
On the only surviving selected-\(T_+\) arcs, the universal concavity theorem
supplies two certified affine deposits.  Starting from \(y_0=H\), set
\[
 y_1^{\rm aff}=y_0+(1-4\alpha)(y_0-\alpha),
\]
\[
 y_2^{\rm aff}
 =y_1^{\rm aff}+(1-5m_3)(y_1^{\rm aff}-m_3).
\]
Whenever an affine coefficient would produce a negative deposit, extensivity
permits the identity wire instead.  The exact CE1 scalar estimate proves that
the realized charge after these two slots crosses the final threshold:
\[
 y_2>e(\delta).
\]
The \(\delta\)-reservoir then returns at least
\[
 1-e(\delta)>1-X,
\]
proving \eqref{eq:body-three-discharge-chain}.  The mechanism is therefore
\[
 \boxed{\text{deposit at }\alpha
 \ \longrightarrow\ \text{deposit at }m_3
 \ \longrightarrow\ \text{trigger at }\delta.}
\]

\paragraph{The \(\CEtwo\) one-hit rule.}
The signed CE2 inequalities prove
\[
 X>e(\alpha),
 \qquad
 \min\{e(\alpha),e(\delta)\}<H.
\]
If \(e(\alpha)<H\), transport the initial charge unchanged to the
\(\alpha\)-slot; that reservoir raises it above \(1-H\).  Otherwise
\[
 e(\delta)<H\le e(\alpha)<X,
\]
so transport the initial charge to the \(\delta\)-slot; that reservoir raises
it above \(1-H\).  Every unused slot is an identity wire.  Thus one of the two
diagonals must fire, although both may fire.

\subsection{The remaining Strategy~2 branches}

\paragraph{The \(N_+=0\) T3-like endpoint branch.}
After support isolation, the exact four-label audit proves
\[
 F_{C_1}(A_1)+F_{C_5}(A_5)<1.
\]
By \eqref{eq:body-endpoint-release-equivalence}, the two endpoint diagonals
release more than the endpoint tax.  The three ordinary internal V triangles
cannot absorb the positive remainder.  The hard four-label calculation is
retained; discharging clarifies its global role but does not replace it by an
unsupported common envelope.

\paragraph{The T3-like and adjacent-\(\Vdone\) rescuers.}
The local rescuer profile forces an opposite endpoint demand
\[
 h\ge g_c^{\rm sc},
\]
while the neighboring strict-supercritical V triangle sends out
\[
 b<g_c^{\rm sc}.
\]
The surplus \(h-b>0\) is discharged through four ordinary
nonsupercritical V triangles.  Their required total boundary contribution is
\(4+(h-b)>4\), exceeding their total capacity four.

\paragraph{Adjacent and nonadjacent Vd placements.}
The residual boundary charges are first transported along ordinary edges.  In
the adjacent placement, the charge \(H\) yields a quarter radial margin and
the center exit satisfies \(\delta<H/4\).  In a nonadjacent placement, the
available charge is \(\min\{A,H\}\), while the center radial reach is strictly
smaller.  These are boundary-to-radial discharge comparisons, but their
conversion factors come from the Vd corner geometry and are not consequences
of the ordinary \(g\)-map alone.

\paragraph{The Vd1--supercritical replacement.}
The exceptional adjacent pair is replaced in two distinct vertex charts by
open nonsupercritical Vd0 roles preserving the full skeleton.  The replacement
step is genuinely geometric: a scalar charge does not encode the required
orientation, adjacent-arm exclusion, and simultaneous radial coverage.  Once
the replacement is made, the all-Vd0 discharging obstruction applies.

\subsection{Limits of the universal discharging language}

A positive charge depending only on the diagonal demand cannot exist.  For
each \(c\), the exact capped map has a Full state with
\[
 F_c(a)=1-a,
 \qquad
 \mathfrak d_c(a)=0.
\]
Consequently neither a fixed rule \(q(c)>0\) nor a change of potential with a
uniform positive increment can hold on the full exact domain.  The charge must
depend on the arriving boundary state.

Likewise, the length of a V-gap is not sufficient state data.  A singleton
V-gap has zero length but remains active because both incident roles are open.
The residual-component operator and the active-gap convention must therefore
be retained.

Finally, the \(N_+=1\), rank-zero all-Vd0 state is not closed by boundary
capacity discharging.  One negative supercritical slack can balance the five
nonnegative slacks.  Its active proof is the center-independent nine-point
obstruction, which uses interior witnesses rather than only skeleton charge.

\begin{proposition}[Branches closed by discharging]
\label{prop:body-discharging-branches}
The diagonal-discharge framework closes all Strategy~2 branches previously
closed by relaxed \(g\)-chains:
\begin{enumerate}
\item the \(N_+=0\) all-\(\Vdzero\) branches of active-gap rank
      \(0,1,2\);
\item the \(N_+=0\) T3-like branch with no Vd1/Vd2 role;
\item every nonempty-gap \(N_+=1\) all-\(\Vdzero\) state;
\item the exactly-one-T3-like and adjacent-rescuer states;
\item the residual and radial-separation terminals in the one-Vd1/Vd2 branch,
      followed, where necessary, by the two-chart replacement.
\end{enumerate}
The exact \(g\)-maps, endpoint label audits, and Vd profiles remain the local
proofs of the stated release and conversion inequalities.
\end{proposition}

\begin{proof}
Items 1 and 2 use the cyclic capacity sum or the paired endpoint release
\eqref{eq:body-endpoint-release-equivalence} followed by the path budget.
Item 3 is the CE1 deposit cascade or CE2 one-hit threshold rule.  Item 4 is the
strict source--sink surplus \(h-b>0\).  Item 5 uses the exact residual and Vd
radial margins, with the all-Vd0 terminal after replacement.  Every local
release is connected to an actual V triangle by
Proposition~\ref{prop:body-exact-discharging}.
\end{proof}
